\section*{第 \ref{ch:b1:central-forces} 章——中心力：行星与卫星}

\begin{solution}{exo:b1:central-forces:1}
$r = \qty{6770}{km}$: $v = \sqrt{GM/r} = \qty{7.7}{km/s}$; $T = 2\pi r/v
= \qty{5.5e3}{s} = \qty{92}{min}$；$1440/92 \approx 16$，即每天 16 次日出。
\end{solution}

\begin{solution}{exo:b1:central-forces:2}
$r = (GMT^2/4\pi^2)^{1/3} = (3.99\times10^{14} \times 7.42\times10^9/39.5)^{1/3}
= \qty{4.22e7}{m}$：高度 \qty{35800}{km}；$v = 2\pi r/T = \qty{3.1}{km/s}$。
\end{solution}

\begin{solution}{exo:b1:central-forces:3}
$v_e = \sqrt{2GM/R}$：地球 \qty{11.2}{km/s}；月球 $\sqrt{2 \times 4.90\times
10^{12}/1.74\times10^6} = \qty{2.4}{km/s}$。气体分子在几百
米每秒的分子，其速率分布在
\qty{2.4}{km/s} 以上仍有尾巴；经过地质时间，月球上的气体已逸散殆尽。
\end{solution}

\begin{solution}{exo:b1:central-forces:4}
$T = a^{3/2}$（年，天文单位）：火星 \qty{1.88}{yr}；木星 \qty{11.9}{yr}；
\qty{30}{AU}：\qty{164}{yr}。
\end{solution}

\begin{solution}{exo:b1:central-forces:5}
$\Delta E = \tfrac12 mv^2 + GMm(1/R - 1/r) = 2.94\times10^{10} + 3.99\times
10^{17} \times (1.570 - 1.477)\times10^{-7} = 2.94\times10^{10} + 0.37\times
10^{10} = \qty{3.3e10}{J}$——\qty{33}{MJ/kg}，即送入轨道每千克所需的能量，相当于 \qty{0.7}{kg}
汽油的能量（火箭所需远多于此，
因为它们还须托起自身的燃料）。
\end{solution}

\begin{solution}{exo:b1:central-forces:6}
$a = (0.586 + 35.1)/2 = \qty{17.8}{AU}$; $e = (35.1 - 0.586)/35.7 = 0.967$;
$T = 17.8^{3/2} = \qty{75}{yr}$。$v_p^2 = GM(2/r_p - 1/a)$，其中 $r_p =
8.77\times10^{10}$ m, $a = 2.67\times10^{12}$ m: $v_p = \qty{54}{km/s}$;
$v_a = v_pr_p/r_a = \qty{0.91}{km/s}$.
\end{solution}

\begin{solution}{exo:b1:central-forces:7}
$a = \qty{24490}{km}$。$v_1 = \qty{7.67}{km/s}$；在 $r_1$ 处的椭圆轨道上：
$\sqrt{GM(2/r_1 - 1/a)} = \qty{10.07}{km/s}$: $\Delta v_1 = \qty{2.4}{km/s}$.
在 $r_2$ 处：$\sqrt{GM(2/r_2 - 1/a)} = \qty{1.62}{km/s}$，圆轨道速率
\qty{3.07}{km/s}，$\Delta v_2 = \qty{1.45}{km/s}$。总计 \qty{3.85}{km/s}；
时长 $\pi\sqrt{a^3/GM} = \qty{1.9e4}{s} = \qty{5.3}{h}$。
\end{solution}

\begin{solution}{exo:b1:central-forces:8}
$E_{p,\mathrm{eff}}' = GMm/r^2 - L^2/mr^3 = 0$，得 $r_c = L^2/GMm^2$。
$E_{p,\mathrm{eff}}'' = -2GMm/r^3 + 3L^2/mr^4$；在 $r_c$ 处，$L^2 = GMm^2r_c$
得 $E'' = GMm/r_c^3$，故 $\omega_r^2 = E''/m = GM/r_c^3 = \omega_{\text{轨道}}^2$。
小的径向扰动恰好每转一周振荡一次：轨道
自行闭合——是一个略微偏心的椭圆，而非玫瑰线。
\end{solution}

\begin{solution}{exo:b1:central-forces:9}
只有引力同时作用于空间站与宇航员，二者沿同一轨道下落，
因此二者之间无需接触力——没有地板向上推，
因此称为“失重”。$g = GM/r^2 = 3.99\times10^{14}/(6.77\times10^6)^2
= \qty{8.7}{m/s^2}$：引力几乎达到全强度；重力并未消失，
它全部用于使运动轨迹弯曲。
\end{solution}

\begin{solution}{exo:b1:central-forces:10}
$a = 4\pi^2r/T^2 = 39.5 \times 3.84\times10^8/(2.36\times10^6)^2 =
\qty{2.7e-3}{m/s^2}$；$g(R/r)^2 = 9.81/60.3^2 = \qty{2.7e-3}{m/s^2}$：
使苹果下落的同一引力定律也维系着月球，只是按 $1/r^2$ 衰减。
\end{solution}

\begin{solution}{exo:b1:central-forces:11}
$E_k = 2Ze^2/4\pi\varepsilon_0r_{\min}$: $r_{\min} = 8.99\times10^9 \times 2 \times 79
\times (1.60\times10^{-19})^2/8.0\times10^{-13} = \qty{4.5e-14}{m} = \qty{45}{fm}$,
为原子核半径的六倍：$\alpha$ 粒子在接触前即折返。
若要达到 \qty{7}{fm}：$E_k \approx \qty{32}{MeV}$。
\end{solution}

\begin{solution}{exo:b1:central-forces:12}
$E_m = -GMm/2r$ 减小 $\Rightarrow$ $r$ 减小 $\Rightarrow$ $v =
\sqrt{GM/r}$ 增大。动能 $GMm/2r$ 增加 $\abs{\Delta E}$，
而势能 $-GMm/r$ 减少 $2\abs{\Delta E}$：所释放的势能有一半
转化为动能，另一半
因阻力而以热量形式耗散。
\end{solution}

\begin{solution}{pb:b1:central-forces:1}
\textbf{1.} $T = \qty{43082}{s}$: $r = (GMT^2/4\pi^2)^{1/3} = (3.986\times10^{14}
\times 1.856\times10^9/39.48)^{1/3} = \qty{2.656e7}{m}$；高度
\qty{20190}{km}。

\textbf{2.} $v = \sqrt{GM/r} = \qty{3.87}{km/s}$.

\textbf{3.} $a = v^2/r = GM/r^2 = \qty{0.565}{m/s^2} = g(R/r)^2$，比地表弱
十七倍——但并非为零：卫星
绕地球持续下落。

\textbf{4.} $E_m/m = -GM/2r = \qty{-7.5}{MJ/kg}$; $E_k/m = \qty{+7.5}{MJ/kg}$,
$E_p/m = \qty{-15.0}{MJ/kg}$.

\textbf{5.} 相对于恒星，地球每个恒星日自转一周
（即轨道参考系）；两圈之后卫星回到
同一地面点的上方：地面轨迹每个恒星日重复一次。

\textbf{6.} $\sin\beta = R/r = 0.240$：$\beta = 13.9^\circ$，一个
宽 $28^\circ$ 的圆盘。

\textbf{7.} 地球静止轨道 \qty{42200}{km}：高 $1.59$ 倍；
导航轨道低于它。

\textbf{8.} $GM(1/R - 1/2r) = 3.986\times10^{14}(1.570\times10^{-7} - 1.88\times
10^{-8}) = \qty{5.5e7}{J/kg} = \qty{55}{MJ/kg}$.

\textbf{9.} $r_1 = \qty{6771}{km}$: $v = \qty{7.67}{km/s}$.

\textbf{10.} $a = \qty{16665}{km}$；在 $r_1$ 处：$\sqrt{GM(2/r_1 - 1/a)} =
\qty{9.69}{km/s}$，$\Delta v_1 = \qty{2.0}{km/s}$；在 $r_2$ 处：\qty{2.47}{km/s}
而圆轨道速率为 \qty{3.87}{km/s}，$\Delta v_2 = \qty{1.4}{km/s}$；合计
\qty{3.4}{km/s}.

\textbf{11.} $\pi\sqrt{a^3/GM} = \qty{1.07e4}{s} \approx \qty{3.0}{h}$.

\textbf{12.} 从低轨道逃逸：$(\sqrt2 - 1)\times 7.67 = \qty{3.2}{km/s}$：
任务的 \qty{3.4}{km/s} 超过了这一数值——到达该轨道的代价
几乎与彻底脱离地球相当。

\textbf{13.} $m_0/m_1 = \eu^{3.4/3.0} = 3.1$: $m_1 = \qty{640}{kg}$,
推进剂 \qty{1360}{kg}——约占总质量的三分之二。

\textbf{14.} $\cos\theta = 0.240$，$\theta = 76^\circ$；比例为 $(1 - 0.240)/2
= 0.38$.

\textbf{15.} 平均而言 $24 \times 0.38 = 9$，
几乎总是超过四颗。

\textbf{16.} 天顶：$2.02\times10^7/c = \qty{67}{ms}$；地平线：距离
$\sqrt{r^2 - R^2} = \qty{2.58e7}{m}$, \qty{86}{ms}.

\textbf{17.} $3/c = \qty{10}{ns}$.

\textbf{18.} 接收机钟差是除其三个坐标外的第四个未知量，
因此需要四个距离、四个方程。

\textbf{19.} $10^{-8}\ \unit{s}/\qty{86400}{s} \approx \num{1e-13}$.

\textbf{20.} $v^2/2c^2 = (3874)^2/(2 \times 8.99\times10^{16}) = \num{8.3e-11}$:
$8.3\times10^{-11} \times 86400 = \qty{7.2}{\micro s}$，即每天走慢。

\textbf{21.} $GM(1/R - 1/r)/c^2 = 3.986\times10^{14} \times 1.19\times10^{-7}/
8.99\times10^{16} = \num{5.3e-10}$：每天 \qty{45.7}{\micro s}，走快。

\textbf{22.} 净值 $+\qty{38.5}{\micro s}$ 每天；$c \times 38.5\times10^{-6}
= \qty{11.5}{km}$，即每天产生的误差。

\textbf{23.} $(10.23 - 10.22999999543)/10.23 = \num{4.47e-10}$，与
净值 $4.5\times10^{-10}$ 相符（\qty{86400}{s} 内为 \qty{38.5}{\micro s}）。

\textbf{24.} 地球赤道隆起、月球与太阳的
吸引、太阳辐射压力（以及小推力器点火）：地面段
连续测量轨道并上传新的星历表，
再由卫星播发。

\textbf{25.} 能量与角动量可由周期确定一条圆轨道；
开普勒第三定律确定其位置；几何关系决定覆盖范围与光行时间；
而时钟——系统的核心——必须做相对论修正，
每天 \qty{38}{\micro s}，否则到夜幕降临时
定位将漂移十千米。
\end{solution}
